\section{Całka Riemanna}
\label{ma1:ch:riemann-integral}
% Import: Analiza matematyczna I, rozdział 29; Teoria miary i całki, rozdział 08.
\begin{definicja}
\label{ma1:def:riemann-sum}
Dla podziału $P$ i punktów wyróżnionych $\xi_i\in[x_{i-1},x_i]$ definiujemy sumę Riemanna
\[
 S(f,P,\xi)=\sum_{i=1}^n f(\xi_i)\Delta x_i.
\]
Granica sum Riemanna przy $|P|\to0$ wynosi $A$, jeżeli dla każdego $\varepsilon>0$ istnieje $\delta>0$ takie, że dla każdego podziału $P$ o $|P|<\delta$ i każdego wyboru $\xi_i$ mamy $|S(f,P,\xi)-A|<\varepsilon$.
\end{definicja}
Dla funkcji ograniczonej zawsze $L(f,P)\le S(f,P,\xi)\le U(f,P)$, ponieważ $m_i\le f(\xi_i)\le M_i$.

\begin{stwierdzenie}
Każda funkcja całkowalna w sensie Riemanna jest ograniczona.
\end{stwierdzenie}
\begin{proof}
Gdyby $f$ była nieograniczona, dla dowolnego ustalonego podziału musiałaby być nieograniczona na co najmniej jednej komórce. Zmieniając znacznik tylko na niej, uzyskalibyśmy sumy o dowolnie dużej różnicy, wbrew definicji całki.
\end{proof}

\begin{twierdzenie}
\label{ma1:thm:riemann-darboux-equivalence}
Dla ograniczonej $f$ całkowalność z definicji \ref{ma1:def:riemann-integral} jest równoważna istnieniu skończonej granicy sum Riemanna przy $|P|\to0$. W obu przypadkach otrzymujemy tę samą wartość.
\end{twierdzenie}
\begin{proof}
Załóżmy całkowalność i oznaczmy całkę przez $A$. Ustalmy $\varepsilon>0$ i podział $P_0$ z $U(f,P_0)-L(f,P_0)<\varepsilon/2$. Niech $k$ oznacza liczbę jego punktów wewnętrznych, a $B>0$ będzie ograniczeniem $|f|$.

Dla dowolnego podziału $P$ weźmy wspólne rozdrobnienie $Q$ podziałów $P,P_0$. Podprzedział $P$ zmienia się w tym rozdrobnieniu tylko wtedy, gdy zawiera w swoim wnętrzu punkt $P_0$. Takich podprzedziałów jest co najwyżej $k$, ponieważ ich wnętrza są parami rozłączne i w każdym znajduje się przynajmniej jeden punkt wewnętrzny $P_0$. Ich łączna długość jest nie większa niż $k|P|$. Punkty $P_0$ pokrywające się z końcami podprzedziałów $P$ nie powodują ich rozdrobnienia. Na każdym z nich zarówno stary wkład do sumy górnej, jak i suma nowych wkładów należą do przedziału $[-B\Delta x_i,B\Delta x_i]$. Różnica jest więc nie większa niż $2B\Delta x_i$. Analogicznie dla sumy dolnej. Z lematu \ref{ma1:lem:darboux-refinement} wynika
\[
 \begin{split}
 U(f,P)&\le U(f,Q)+2Bk|P|\le U(f,P_0)+2Bk|P|,\\
 L(f,P)&\ge L(f,Q)-2Bk|P|\ge L(f,P_0)-2Bk|P|.
 \end{split}
\]
Ponieważ $L(f,P_0)\le A\le U(f,P_0)$ oraz $L(f,P)\le S(f,P,\xi)\le U(f,P)$, mamy
\[
 -\bigl(U(f,P_0)-L(f,P_0)\bigr)-2Bk|P|
 \le S(f,P,\xi)-A
 \le U(f,P_0)-L(f,P_0)+2Bk|P|.
\]
Każda suma Riemanna dla $P$ spełnia zatem
\[
 |S(f,P,\xi)-A|<\varepsilon/2+2Bk|P|.
\]
Wybieramy $\delta=\varepsilon/(4B(k+1))$. Wtedy prawa strona jest mniejsza od $\varepsilon$ dla $|P|<\delta$, niezależnie od wyboru punktów wyróżnionych.

Odwrotnie, załóżmy warunek graniczny dla sum Riemanna z granicą $A$. Dla $\varepsilon>0$ wybieramy $\delta$ odpowiadające dokładności $\varepsilon/4$ i podział $P$ o średnicy mniejszej od $\delta$ (na przykład podział na $n$ równych części dla $n>(b-a)/\delta$). Z definicji kresów dla każdego $i$ można wybrać $\xi_i^+,\xi_i^-\in I_i$ tak, aby
\[
 M_i-f(\xi_i^+)<\frac{\varepsilon}{4(b-a)},\qquad
 f(\xi_i^-)-m_i<\frac{\varepsilon}{4(b-a)}.
\]
Po pomnożeniu przez długości i zsumowaniu otrzymujemy
$U(f,P)<S(f,P,\xi^+)+\varepsilon/4<A+\varepsilon/2$
oraz $L(f,P)>S(f,P,\xi^-)-\varepsilon/4>A-\varepsilon/2$.
Zatem $U(f,P)-L(f,P)<\varepsilon$. Kryterium Darboux daje całkowalność. Jeżeli $J$ oznacza uzyskaną całkę, to $L(f,P)\le J\le U(f,P)$, więc $|J-A|<\varepsilon/2$. Ponieważ $\varepsilon>0$ było dowolne, mamy $J=A$.
\end{proof}
W zapisie ze wstępu $S(f,P,\xi)=S(f,\mathcal P)$, gdzie $\mathcal P$ jest podziałem znaczonym, a $\xi_i$ są jego znacznikami. Istotna jest dowolność podziałów i punktów wyróżnionych. Zbieżność sum dla jednego szczególnego ciągu podziałów i jednego wyboru punktów nie wystarcza.

\subsection{Kryterium Lebesgue'a}
\begin{twierdzenie}[Kryterium Lebesgue'a całkowalności Riemanna]
\label{thm:riemann-criterion}
Niech $f:[a,b]\to\mathbb R$ będzie ograniczona, gdzie
$a<b$. Wówczas $f$ jest całkowalna w sensie Riemanna
tedy i tylko wtedy, gdy zbiór punktów nieciągłości $f$
ma miarę Lebesgue'a zero.
\end{twierdzenie}
\begin{proof}
Dla $x\in[a,b]$ oznaczmy przez $\omega_f(x)$
oscylację $f$ w punkcie, czyli
\[
\omega_f(x):=\inf_{r>0}
\left(\sup_{y\in[a,b]\cap(x-r,x+r)}f(y)
-\inf_{y\in[a,b]\cap(x-r,x+r)}f(y)\right).
\]
Funkcja $f$ jest ciągła w $x$ wtedy i tylko wtedy,
gdy $\omega_f(x)=0$. Dla $\varepsilon>0$ zbiór
$D_\varepsilon=\{x:\omega_f(x)\ge\varepsilon\}$
jest domknięty: jeśli $\omega_f(x)<\varepsilon$,
to na pewnym otoczeniu punktu $x$ oscylacja $f$
jest mniejsza niż $\varepsilon$, a zatem także
oscylacje w punktach pobliskich są mniejsze od
$\varepsilon$. Zatem $D_\varepsilon$ jest zwarty.

Załóżmy najpierw całkowalność Riemanna i
ustalmy $\varepsilon>0$. Weźmy dowolny podział
$P=(a=x_0<\cdots<x_N=b)$. Jeśli
$x\in D_\varepsilon$ nie jest punktem podziału,
to oscylacja $f$ na domkniętym przedziale
$[x_{j-1},x_j]$ zawierającym $x$ jest co najmniej
$\varepsilon$. Dlatego $D_\varepsilon$ mieści się
w sumie skończenie wielu punktów podziału i tych
przedziałów, których oscylacja jest nie mniejsza
niż $\varepsilon$. Suma długości tych przedziałów
jest nie większa niż
$\varepsilon^{-1}(U(f,P)-L(f,P))$.
Ponieważ dla funkcji całkowalnej w sensie Riemanna
różnicę tę można uczynić dowolnie małą,
$\lambda_1(D_\varepsilon)=0$.
Zbiór nieciągłości jest sumą
$\bigcup_{m\ge1}D_{1/m}$, więc jest zerowy.

Odwrotnie, niech zbiór nieciągłości $D$ ma miarę
zero, a $|f|\le M$. Ustalmy $\varepsilon>0$.
Zbiór $D_\varepsilon$ jest zwarty i zerowy, więc
istnieje zbiór otwarty $O$ zawierający $D_\varepsilon$
o dowolnie małej mierze Lebesgue'a (regularność zewnętrzna miary Lebesgue'a). Wybieramy przy tym
$\lambda_1(O)<b-a$, dzięki czemu zwarty zbiór
$K=[a,b]\setminus O$ jest niepusty. Dla każdego
$x\in K$ mamy $\omega_f(x)<\varepsilon$,
zatem istnieje promień $r_x>0$, na którego kuli
$B(x,r_x)\cap[a,b]$ oscylacja $f$ jest
mniejsza niż $\varepsilon$. Zwartość $K$ pozwala
wybrać skończenie wiele kul $B(x_j,r_{x_j}/2)$
pokrywających $K$. Połóżmy
$\delta=\frac12\min_j r_{x_j}>0$.
Jeżeli domknięty przedział $J\subseteq[a,b]$
o długości mniejszej niż $\delta$ przecina $K$,
wybieramy $y\in J\cap K$ i indeks $j$ taki, że
$y\in B(x_j,r_{x_j}/2)$. Dla każdego $z\in J$
mamy
$|z-x_j|\le|z-y|+|y-x_j|<\delta+r_{x_j}/2
\le r_{x_j}$.
Zatem $J\subseteq B(x_j,r_{x_j})$, gdzie
oscylacja $f$ jest mniejsza niż $\varepsilon$.

Weźmy podział $P$ o średnicy mniejszej niż $\delta$.
Na przedziałach przecinających $K$ oscylacja $f$
jest mniejsza niż $\varepsilon$. Pozostałe
przedziały (z wyjątkiem ewentualnych końców)
mieszczą się w $O$, więc suma ich długości
nie przekracza $\lambda_1(O)$. Otrzymujemy
\[
U(f,P)-L(f,P)
\le\varepsilon(b-a)+2M\lambda_1(O).
\]
Wybierając najpierw $\varepsilon$, a potem
$\lambda_1(O)$ dostatecznie małe, czynimy prawą
stronę dowolnie małą. Kryterium Darboux dowodzi
całkowalności Riemanna.
\end{proof}

\subsection{Własności całki oznaczonej}
\begin{twierdzenie}
\label{ma1:thm:riemann-algebra}
Jeżeli $f,g$ są całkowalne na $[a,b]$ oraz $\alpha,\beta\in\R$, to $\alpha f+\beta g$, $fg$ oraz $|f|$ są całkowalne. Ponadto
\[
 \int_a^b(\alpha f+\beta g)\,\dd x
 =\alpha\int_a^b f\,\dd x+\beta\int_a^b g\,\dd x.
\]
\end{twierdzenie}
\begin{proof}
Dla każdego podziału z punktami wyróżnionymi suma Riemanna funkcji $\alpha f+\beta g$ jest sumą odpowiednich sum pomnożonych przez skalary. Jeżeli $A,C$ oznaczają całki $f,g$, to dla zadanej $\varepsilon>0$ wybieramy wspólną średnicę graniczną dla dokładności $\varepsilon/(1+|\alpha|+|\beta|)$ obu sum. Nierówność trójkąta daje błąd mniejszy od $\varepsilon$ względem $\alpha A+\beta C$. Twierdzenie \ref{ma1:thm:riemann-darboux-equivalence} dowodzi całkowalności i wzoru.

Niech $|f|,|g|\le B$ dla $B>0$. Dla dowolnych $s,t$ w jednym podprzedziale
\[
 |f(s)g(s)-f(t)g(t)|\le B|g(s)-g(t)|+B|f(s)-f(t)|.
\]
Oscylacja ograniczonej funkcji $h$ na niepustym zbiorze $E$ jest równa $\sup_{s,t\in E}|h(s)-h(t)|$. Istotnie, dla $m=\inf_E h$ i $M=\sup_E h$ zachodzi $|h(s)-h(t)|\le M-m$. Dla dowolnego $\eta>0$ wybieramy $s,t\in E$ z $h(s)>M-\eta/2$ i $h(t)<m+\eta/2$. Wówczas $|h(s)-h(t)|\ge h(s)-h(t)>M-m-\eta$. Dowolność $\eta$ daje przeciwną nierówność. Zatem oscylacja iloczynu na podprzedziale jest nie większa niż $B$ razy suma oscylacji czynników. Wybieramy wspólne rozdrobnienie podziałów, dla których różnice sum $f$ i $g$ są mniejsze od $\varepsilon/(2B)$. Na wspólnym rozdrobnieniu $Q$ różnice nadal spełniają te oszacowania, zatem
\[
 U(fg,Q)-L(fg,Q)
 \le B\bigl(U(f,Q)-L(f,Q)+U(g,Q)-L(g,Q)\bigr)
 <\varepsilon.
\]
Iloczyn jest ograniczony przez $B^2$, więc stosujemy kryterium Darboux. Dla modułu korzystamy z $||f(s)|-|f(t)||\le |f(s)-f(t)|$, co daje nie większą oscylację niż dla $f$. Kryterium Darboux kończy dowód.
\end{proof}

\begin{twierdzenie}
\label{ma1:thm:riemann-order}
Dla całkowalnych $f,g$ na $[a,b]$ z $f\le g$ wynika $\int_a^b f\,\dd x\le\int_a^b g\,\dd x$. Jeżeli $m\le f\le M$, to
\[
 m(b-a)\le\int_a^b f\,\dd x\le M(b-a),\qquad
 \left|\int_a^b f\,\dd x\right|\le\int_a^b |f|\,\dd x.
\]
\end{twierdzenie}
\begin{proof}
Dla każdego podziału $S(f,P,\xi)\le S(g,P,\xi)$, więc przejście do granicy zachowuje nieostrą nierówność. Całka funkcji stałej $c$ wynosi $c(b-a)$, gdyż każda jej suma Riemanna jest równa tej liczbie. Porównujemy $f$ z funkcjami stałymi $m,M$. Nierówności $-|f|\le f\le |f|$ wraz z liniowością i całkowalnością modułu dają ostatnią tezę.
\end{proof}

\begin{twierdzenie}[Addytywność względem przedziału]
\label{ma1:thm:riemann-additivity}
Niech $a<c<b$ i niech $f$ będzie ograniczona na $[a,b]$. Jest ona całkowalna na $[a,b]$ wtedy i tylko wtedy, gdy jej ograniczenia do $[a,c]$ i $[c,b]$ są całkowalne. Wówczas
\[
 \int_a^b f\,\dd x=\int_a^c f\,\dd x+\int_c^b f\,\dd x.
\]
\end{twierdzenie}
\begin{proof}
Jeżeli $f$ jest całkowalna na całym przedziale, dla $\varepsilon>0$ wybieramy podział z różnicą sum mniejszą od $\varepsilon$ i dodajemy do niego punkt $c$. Różnica nie wzrasta. Jest ona sumą dwóch nieujemnych różnic dla podziałów ograniczonych do obu przedziałów, więc każda z nich jest mniejsza od $\varepsilon$. Kryterium Darboux daje całkowalność ograniczeń.

Odwrotnie, dla obu ograniczeń wybieramy podziały z różnicą sum mniejszą od $\varepsilon/2$ i łączymy je. Różnica sum na całym przedziale jest sumą obu różnic, więc jest mniejsza od $\varepsilon$. Aby wyznaczyć wartość całki, oznaczmy całki na obu częściach przez $A_1,A_2$. Dla połączonego podziału $P$ liczba $A_1+A_2$ leży pomiędzy $L(f,P)$ i $U(f,P)$, podobnie jak całka na całym przedziale. Moduł różnicy tych dwóch liczb jest mniejszy od $\varepsilon$. Dowolność $\varepsilon$ daje wzór.
\end{proof}
Przyjmujemy $\int_a^a f\,\dd x=0$ oraz $\int_b^a f\,\dd x=-\int_a^b f\,\dd x$ dla $a<b$. Wówczas addytywność zachodzi dla dowolnego porządku trzech końców należących do przedziału całkowalności. Istotnie, dla trzech różnych punktów $r<s<t$ oznaczmy $A=\int_r^s f\,\dd x$ i $C=\int_s^t f\,\dd x$. Wówczas całki w kierunku rosnącym mają wartości $A,C,A+C$, a w kierunku malejącym ich przeciwności. Przypisując punktom $r,s,t$ odpowiednio liczby $0,A,A+C$, każdą całkę między tymi punktami otrzymujemy jako różnicę liczby przypisanej końcowi i początkowi. Takie różnice spełniają addytywność w dowolnej kolejności. Jeżeli dwa punkty są równe, teza wynika z całki zerowej lub zmiany orientacji. Z twierdzenia \ref{ma1:thm:riemann-additivity} wynika też całkowalność na każdym domkniętym podprzedziale.

\begin{wniosek}
\label{ma1:cor:riemann-finite-changes}
Zmiana wartości funkcji całkowalnej w skończonej liczbie punktów nie zmienia jej całkowalności ani wartości całki.
\end{wniosek}
\begin{proof}
Niech $g$ różni się od $f$ jedynie w $k$ punktach. Różnica $h=g-f$ jest ograniczona, powiedzmy $|h|\le B$, oraz zerowa poza tymi punktami. Dla dowolnego podziału co najwyżej $2k$ jego podprzedziałów zawiera któryś z tych punktów (punkt może być wspólnym końcem dwóch). Zatem $|S(h,P,\xi)|\le 2kB|P|$. Sumy dążą do zera przy $|P|\to0$, więc $h$ jest całkowalna i jej całka wynosi zero. Liniowość dla $g=f+h$ daje tezę.
\end{proof}

\subsection{Twierdzenie o wartości średniej i przykłady}
\begin{twierdzenie}[O wartości średniej dla całek]
\label{ma1:thm:riemann-mean-value}
Niech $f$ będzie ciągła na $[a,b]$ i niech $g$ będzie całkowalna oraz nieujemna. Istnieje $\xi\in[a,b]$ takie, że
\[
 \int_a^b fg\,\dd x=f(\xi)\int_a^b g\,\dd x.
\]
W szczególności $\int_a^b f\,\dd x=f(\xi)(b-a)$ dla pewnego $\xi\in[a,b]$.
\end{twierdzenie}
\begin{proof}
Niech $m,M$ będą minimum i maksimum $f$, istniejące z twierdzenia Weierstrassa. Z $g\ge0$ wynika $mg\le fg\le Mg$. Całkowalność iloczynu i porównanie całek dają
$m\int g\le\int fg\le M\int g$, gdzie wszystkie całki są po $[a,b]$.
Jeżeli $\int g=0$, to również $\int fg=0$, więc dowolne $\xi$ spełnia wzór. Jeżeli $\int g>0$, dzielimy przez tę liczbę i otrzymujemy iloraz leżący w $[m,M]$. Twierdzenie o wartości pośredniej daje punkt $\xi$ o tej wartości $f$. Wybór $g=1$ daje szczególny przypadek.
\end{proof}
\begin{przyklad}
\label{ma1:ex:riemann-square}
Obliczymy $\int_0^1 x^2\,\dd x$ bez użycia funkcji pierwotnej. Dla podziału na $n$ równych części i prawych końców jako punktów wyróżnionych otrzymujemy
\[
 S_n=\frac1{n^3}\sum_{i=1}^n i^2
 =\frac{n(n+1)(2n+1)}{6n^3}.
\]
Tożsamość dla sumy kwadratów wynika z indukcji: dla $n=1$ obie strony wynoszą $1$, a krok indukcyjny ma postać
\[
 \frac{n(n+1)(2n+1)}6+(n+1)^2
 =\frac{(n+1)(2n^2+7n+6)}6
 =\frac{(n+1)(n+2)(2n+3)}6.
\] Ponieważ $|P_n|=1/n\to0$ i $x^2$ jest ciągła, twierdzenie \ref{ma1:thm:riemann-darboux-equivalence} daje całkę równą $\lim S_n=1/3$.
\end{przyklad}
\begin{przyklad}
\label{ma1:ex:riemann-step}
Dla $a<c<b$ funkcja $\chi_{[c,b]}$ na $[a,b]$ jest całkowalna (ma tylko jedną nieciągłość). Na $[a,c]$ różni się od zera wyłącznie w $c$, więc jej całka na tym przedziale jest zerowa. Na $[c,b]$ jest stała równa $1$, więc całka wynosi $b-c$. Addytywność daje $\int_a^b\chi_{[c,b]}(x)\,\dd x=b-c$.
\end{przyklad}

\subsection{Podstawowe twierdzenie rachunku całkowego}
\begin{twierdzenie}
\label{ma1:thm:riemann-fundamental}
Niech $f$ będzie całkowalna na $[a,b]$ i $x_0\in[a,b]$. Funkcja
\[
 F(x)=\int_{x_0}^x f(t)\,\dd t,\qquad x\in[a,b],
\]
jest ciągła na $[a,b]$. Jeżeli $f$ jest ciągła w $x\in(a,b)$, to $F$ jest różniczkowalna w $x$ i $F'(x)=f(x)$. Przy ciągłości jednostronnej $f$ na końcu analogiczna równość zachodzi dla pochodnej jednostronnej.
\end{twierdzenie}
\begin{proof}
Niech $|f|\le B$ na $[a,b]$. Addytywność i oszacowanie całki dają dla $u<v$
\[
 |F(v)-F(u)|=\left|\int_u^v f(t)\,\dd t\right|\le B(v-u).
\]
A zatem $F$ jest Lipschitzowska (dla $B=0$ jest stała), w szczególności ciągła.

Ustalmy punkt $x$ ciągłości $f$ i przyrost $h\ne0$ taki, że $x+h\in[a,b]$. Addytywność, liniowość i całka funkcji stałej (z uwzględnieniem orientacji dla $h<0$) dają
\[
 \frac{F(x+h)-F(x)}h-f(x)
 =\frac1h\int_x^{x+h}(f(t)-f(x))\,\dd t.
\]
Dla $\varepsilon>0$ z ciągłości $f$ wybieramy $\delta>0$, aby $|f(t)-f(x)|<\varepsilon/2$, gdy $|t-x|<\delta$ i $t\in[a,b]$. Jeżeli $0<|h|<\delta$, każdy $t$ pomiędzy $x$ i $x+h$ spełnia ten warunek. Stosując oszacowanie całki na przedziale z końcami w kolejności rosnącej, otrzymujemy
\[
 \left|\frac{F(x+h)-F(x)}h-f(x)\right|
 \le\frac{\varepsilon}{2|h|}|h|=\frac{\varepsilon}{2}<\varepsilon.
\]
Z definicji pochodnej dostajemy $F'(x)=f(x)$. Przy końcu przedziału używamy jedynie dopuszczalnych przyrostów i ciągłości z odpowiedniej strony.
\end{proof}
\begin{wniosek}
\label{ma1:cor:continuous-primitive-existence}
Każda funkcja ciągła $f$ na niepustym otwartym przedziale $I$ ma funkcję pierwotną. Dla $x_0\in I$ dokładnie jedna z nich spełnia warunek $F(x_0)=0$ i jest dana wzorem
\[
 F(x)=\int_{x_0}^x f(t)\,\dd t.
\]
\end{wniosek}
\begin{proof}
Dla każdego $x\in I$ domknięty odcinek o końcach $x_0,x$ zawiera się w $I$. Całka jest określona przez twierdzenie \ref{ma1:thm:riemann-continuous} i konwencję orientacji. Aby sprawdzić pochodną w ustalonym $x$, wybieramy domknięty przedział zawarty w $I$, zawierający $x_0$ oraz $x$ we wnętrzu. Na tym przedziale stosujemy twierdzenie \ref{ma1:thm:riemann-fundamental}; otrzymujemy $F'(x)=f(x)$. Całka od $x_0$ do $x_0$ wynosi zero. Jednoznaczność wynika z twierdzenia Lagrange'a: różnica dwóch funkcji pierwotnych ma pochodną zerową, jest więc stała na $I$, a warunek normalizacji wymusza, że stała wynosi zero.
\end{proof}


\begin{twierdzenie}[Wzór Newtona--Leibniza]
\label{ma1:thm:newton-leibniz}
Niech $f$ będzie całkowalna na $[a,b]$. Jeżeli $G$ jest ciągła na $[a,b]$, różniczkowalna na $(a,b)$ i $G'=f$ w jego wnętrzu, to
\[
 \int_a^b f(x)\,\dd x=G(b)-G(a).
\]
Nie wymagamy istnienia pochodnych na końcach.
\end{twierdzenie}
\begin{proof}
Na każdym podprzedziale $[x_{i-1},x_i]$ funkcja $G$ jest ciągła i różniczkowalna w jego wnętrzu, więc twierdzenie Lagrange'a daje $\xi_i\in(x_{i-1},x_i)$ takie, że $G(x_i)-G(x_{i-1})=G'(\xi_i)\Delta x_i=f(\xi_i)\Delta x_i$. Sumując, otrzymujemy dokładnie
\[
 \sum_{i=1}^n\bigl(G(x_i)-G(x_{i-1})\bigr)
 =G(b)-G(a)=S(f,P,\xi).
\] Wybieramy podziały na $n$ równych części i odpowiadające im punkty z twierdzenia Lagrange'a. Ich średnice dążą do zera, a całkowalność $f$ i twierdzenie \ref{ma1:thm:riemann-darboux-equivalence} dają granicę sum równą całce. Lewa strona równości dla sum jest niezależna od podziału, więc otrzymujemy tezę.
\end{proof}
Dla funkcji ciągłej można także odjąć od $G$ funkcję $F$ z twierdzenia \ref{ma1:thm:riemann-fundamental}. Różnica ma zerową pochodną we wnętrzu i jest ciągła na końcach, więc jest stała z twierdzenia Lagrange'a. Powyższy dowód obejmuje jednak również całkowalne pochodne nieciągłe.

\begin{przyklad}
\label{ma1:ex:riemann-newton-leibniz}
Z wzorów na funkcje pierwotne i wzoru Newtona--Leibniza otrzymujemy
\[
 \int_0^1 x^2\,\dd x=\left[\frac{x^3}{3}\right]_0^1=\frac13,\qquad
 \int_1^2\frac{\dd x}{x}=[\log x]_1^2=\log2.
\]
Zapis $[G(x)]_a^b$ oznacza $G(b)-G(a)$, a nie iloraz. Obie funkcje podcałkowe są ciągłe na wskazanych przedziałach, więc są całkowalne. Na przedziale przecinającym zero nie wolno stosować drugiego wzoru, gdyż $1/x$ nie jest tam określona i ograniczona. Nie jest to przypadek klasycznej całki Riemanna na przedziale domkniętym.
\end{przyklad}


\begin{wniosek}
Klasa $\mathcal R[a,b]$ funkcji całkowalnych w sensie Riemanna jest kratą wektorową, a granica jednostajna funkcji z tej klasy także do niej należy i jej całka jest granicą całek.
\end{wniosek}
\begin{proof}
Twierdzenie \ref{ma1:thm:riemann-algebra} daje domkniętość na kombinacje liniowe i moduł, a wzory $\max(f,g)=(f+g+|f-g|)/2$ oraz $\min(f,g)=(f+g-|f-g|)/2$ dają domkniętość na maksimum i minimum punktowe.
Jeżeli $f_n\to f$ jednostajnie, to $f$ jest ograniczona, ponieważ dla pewnego $n$ mamy $\|f-f_n\|_\infty<1$. Dla każdego podziału $P$ zachodzi
\[
 U(f,P)-L(f,P)\le U(f_n,P)-L(f_n,P)+2(b-a)\|f-f_n\|_\infty.
\]
Istotnie, na każdym podprzedziale kres górny może wzrosnąć, a kres dolny zmaleć co najwyżej o $\|f-f_n\|_\infty$. Dla zadanej dokładności $\varepsilon>0$ wybieramy najpierw $n$ z $2(b-a)\|f-f_n\|_\infty<\varepsilon/2$, a potem podział z różnicą sum $f_n$ mniejszą od $\varepsilon/2$. Kryterium Darboux daje całkowalność $f$. Z liniowości i twierdzenia \ref{ma1:thm:riemann-order} otrzymujemy
\[
 \left|\int_a^b f_n\,\mathrm dx-\int_a^b f\,\mathrm dx\right|
 \le (b-a)\|f_n-f\|_\infty\longrightarrow0.
\]
\end{proof}

\begin{uwaga}[Słabości całki Riemanna]
Zbieżność punktowa nie zachowuje całkowalności: jeżeli $\Q\cap[0,1]=\{q_1,q_2,\dots\}$, to $\ind{\{q_1,\dots,q_n\}}\to\ind{\Q\cap[0,1]}$. Całka Riemanna nie obejmuje funkcji nieograniczonych ani wszystkich pochodnych (istnieją różniczkowalne funkcje o nieograniczonej pochodnej). Seminorma $\int_a^b|f|$ po utożsamieniu funkcji równych p.w. jest normą, lecz otrzymana przestrzeń nie jest zupełna.
\end{uwaga}
\begin{przyklad}[Dlaczego brakuje zupełności]
Na $[0,1]$ niech $f_n(x)=\min\{n,x^{-1/2}\}$ dla $x>0$ i $f_n(0)=n$. Każde $f_n$ jest ciągłe, a $f_n\uparrow f$ poza zerem, gdzie $f(x)=x^{-1/2}$. Zbieżność zmajoryzowana (lub bezpośrednie całkowanie) daje $\int_0^1|f_n-f|\,\mathrm dx\to0$, zatem $(f_n)$ jest ciągiem Cauchy'ego w normie $L^1$. Nie istnieje jednak ograniczona funkcja całkowalna w sensie Riemanna, która byłaby równa $f$ prawie wszędzie.
\end{przyklad}
